\section{Билет 18. Выпуклый симплексный метод Зангвилла}

\ticketq{Выпуклый симплексный метод Зангвилла.}

\paragraph{Постановка.}
\[
  \varphi(x) \to \min , \qquad Ax = b , \qquad x \ge 0 ,
\]
$\varphi$ выпукла и дифференцируема, $A$ размера $m \times n$,
$\rank A = m$. Разбиение на базис и небазис: $x = (x_B, x_N)$, $A = (B, N)$,
$B$ невырождена, $x_B > 0$ (невырожденный план).

\textbf{Идея.} Обобщение симплекс-метода ЛП на выпуклую нелинейную цель: как и
там, на каждой итерации меняется ровно \emph{одна небазисная} переменная, а
базисные пересчитываются из $Ax = b$.

\paragraph{Приведённый градиент.}
\[
  r_j = \frac{\partial\varphi}{\partial x_j}
      - \nabla_B\T\varphi\; B^{-1}A_j , \qquad j \in N .
\]

\paragraph{Критерий оптимальности (Куна--Таккера).}
\[
  r_j \ge 0 \ \text{при } x_j = 0 ,
  \qquad
  r_j = 0 \ \text{при } x_j > 0 ,
  \qquad j \in N .
\]
Для выпуклой $\varphi$ эти условия достаточны для глобального минимума.

\paragraph{Выбор переменной.}
\[
  u = \min_{j \in N} r_j = r_s ,
  \qquad
  v = \max_{j \in N} x_j r_j = x_t r_t .
\]
\textbf{Останов:} $u \ge 0$ и $v = 0$ — это и есть условия Куна--Таккера
($r_j \ge 0$ всюду и $r_jx_j = 0$).

Иначе: если $\abs{u} \ge v$ — \emph{увеличиваем} $x_s$ ($\sigma = +1$, $q = s$);
если $\abs{u} < v$ — \emph{уменьшаем} $x_t$ ($\sigma = -1$, $q = t$).

\paragraph{Направление.}
\[
  d_q = \sigma , \qquad d_j = 0 \ \ (j \in N,\ j \ne q) ,
  \qquad d_B = -\sigma B^{-1}A_q .
\]

\textbf{Допустимость по равенствам:}
$Ad = Bd_B + A_q d_q = -\sigma A_q + \sigma A_q = 0$ $\Rightarrow$
$A(x + \alpha d) = b$.

\textbf{Спуск.} Так как меняется одна небазисная компонента,
$\nabla\T\varphi\,d = r_q d_q$, и разбор двух случаев:
\begin{itemize}[nosep,left=0pt]
  \item $\abs{u} \ge v$, $\sigma = +1$: тогда $r_q = u < 0$ (иначе был бы
        останов), значит $\nabla\T\varphi\,d = r_q < 0$; увеличение $x_q \ge 0$
        допустимо всегда;
  \item $\abs{u} < v$, $\sigma = -1$: тогда $r_q > 0$ и $x_q > 0$ (иначе
        $v = 0$), значит $\nabla\T\varphi\,d = -r_q < 0$; уменьшать $x_q$ можно
        до нуля.
\end{itemize}
Базисные компоненты строго положительны, поэтому при малом $\alpha$ точка
остаётся допустимой.

\paragraph{Выбор шага.}
\[
  \alpha_{\max} = \min\Big\{ -\frac{x_i}{d_i} \ :\ d_i < 0 \Big\}
  \quad(\text{при } \sigma = -1 \text{ учитывается и } x_q) ,
\]
\[
  \alpha_k = \argmin_{0 \le \alpha \le \alpha_{\max}} \varphi(x_k + \alpha d) .
\]

\paragraph{Алгоритм.}
\begin{csteps}
  \item Задать допустимый базисный план $x^0$ с базисом $B$.
  \item Вычислить $r_j$, $j \in N$; найти $u = r_s$ и $v = x_t r_t$.
  \item $u \ge 0$ и $v = 0$ — конец, $x^k$ оптимальна.
  \item $\abs{u} \ge v$: $\sigma = +1$, $q = s$; иначе $\sigma = -1$, $q = t$.
  \item Построить $d$: $d_q = \sigma$, прочие небазисные нули,
        $d_B = -\sigma B^{-1}A_q$.
  \item Найти $\alpha_{\max}$ и $\alpha_k$; положить
        $x^{k+1} = x^k + \alpha_k d$.
  \item Если на границе базисная переменная обратилась в нуль — она выходит из
        базиса, $x_q$ входит; пересчитать $B$. К~шагу~2.
\end{csteps}

\paragraph{Теорема о сходимости.}
\textit{$\varphi(x^k)$ убывает; предельные точки $\{x^k\}$ — точки
Куна--Таккера, а при выпуклой $\varphi$ — решения задачи.}

\textsc{Схема.} $\varphi(x^k)$ убывает и ограничена снизу $\Rightarrow$
сходится. Число базисов конечно, поэтому можно выделить подпоследовательность
с фиксированным базисом; при фиксированном базисе направление $d(x)$
непрерывно по $x$. Если для предельной точки $\bar{x}$ было бы
$d(\bar{x}) \ne 0$, нашёлся бы шаг $\bar\alpha > 0$ с
$\varphi(\bar{x} + \bar\alpha d(\bar{x})) < \varphi(\bar{x})$, что
противоречит сходимости $\varphi(x^k)$. Значит $d(\bar{x}) = 0$, то есть
выполнены условия Куна--Таккера. $\blacksquare$

\paragraph{Отличие от метода условного градиента.}
Оба метода работают на многограннике $Ax = b$, $x \ge 0$ и используют
приведённый градиент, но вспомогательная задача разная. В методе
Франка--Вульфа (билет~12) решается \emph{задача ЛП на всём} $S$, и $y_k$ —
вершина многогранника, то есть за шаг меняются сразу все координаты. У
Зангвилла задачи ЛП нет вовсе: по правилу $u$/$v$ выбирается \emph{одна}
небазисная переменная, и шаг делается вдоль одного ребра — ровно как в
симплекс-методе ЛП.

\seefull{гл.~3, \S15}
